%% Use the "normalphoto" option if you want a normal photo instead of cropped to a circle
% \documentclass[10pt,a4paper,normalphoto]{altacv}

\documentclass[10pt,a4paper,ragged2e,withhyper]{altacv}
%% AltaCV uses the fontawesome5 and packages.
%% See http://texdoc.net/pkg/fontawesome5 for full list of symbols.

% Change the page layout if you need to
\geometry{left=1.25cm,right=1.25cm,top=1.25cm,bottom=1.25cm,columnsep=1.25cm}

% The paracol package lets you typeset columns of text in parallel
\usepackage{paracol}

% Change the font if you want to, depending on whether
% you're using pdflatex or xelatex/lualatex
% WHEN COMPILING WITH XELATEX PLEASE USE
% xelatex -shell-escape -output-driver="xdvipdfmx -z 0" sample.tex
\ifxetexorluatex
  % If using xelatex or lualatex:
  \setmainfont{Roboto Slab}
  \setsansfont{Lato}
  \renewcommand{\familydefault}{\sfdefault}
\else
  % If using pdflatex:
  \usepackage[rm]{roboto}
  \usepackage[defaultsans]{lato}
  % \usepackage{sourcesanspro}
  \renewcommand{\familydefault}{\sfdefault}
\fi

% Color palette 
\definecolor{MainColor}{HTML}{780845}
\definecolor{SecondColor}{HTML}{804070}
\definecolor{SlateGrey}{HTML}{2E2E2E}
\definecolor{LightGrey}{HTML}{666666}
\definecolor{GoldenEarth}{HTML}{E7D192}
\colorlet{name}{black}
\colorlet{tagline}{SecondColor}
\colorlet{heading}{MainColor}
\colorlet{headingrule}{GoldenEarth}
\colorlet{subheading}{SecondColor}
\colorlet{accent}{SecondColor}
\colorlet{emphasis}{SlateGrey}
\colorlet{body}{LightGrey}

% Change some fonts, if necessary
\renewcommand{\namefont}{\Huge\rmfamily\bfseries}
\renewcommand{\personalinfofont}{\footnotesize}
\renewcommand{\cvsectionfont}{\LARGE\rmfamily\bfseries}
\renewcommand{\cvsubsectionfont}{\large\bfseries}


% Change the bullets for itemize and rating marker
% for \cvskill if you want to
\renewcommand{\cvItemMarker}{{\small\textbullet}}
\renewcommand{\cvRatingMarker}{\faCircle}
% ...and the markers for the date/location for \cvevent
% \renewcommand{\cvDateMarker}{\faCalendar*[regular]}
% \renewcommand{\cvLocationMarker}{\faMapMarker*}

%% Use (and optionally edit if necessary) this .tex if you
%% want to use an author-year reference style like APA(6)
%% for your publication list
% % When using APA6 if you need more author names to be listed
% because you're e.g. the 12th author, add apamaxprtauth=12
\usepackage[backend=biber,style=apa6,sorting=ydnt]{biblatex}
\defbibheading{pubtype}{\cvsubsection{#1}}
\renewcommand{\bibsetup}{\vspace*{-\baselineskip}}
\AtEveryBibitem{%
  \makebox[\bibhang][l]{\itemmarker}%
  \iffieldundef{doi}{}{\clearfield{url}}%
}
\setlength{\bibitemsep}{0.25\baselineskip}
\setlength{\bibhang}{1.25em}


%% Use (and optionally edit if necessary) this .tex if you
%% want an originally numerical reference style like IEEE
%% for your publication list
\usepackage[backend=biber,style=ieee,sorting=ydnt,defernumbers=true]{biblatex}
%% For removing numbering entirely when using a numeric style
\setlength{\bibhang}{1.25em}
\DeclareFieldFormat{labelnumberwidth}{\makebox[\bibhang][l]{\itemmarker}}
\setlength{\biblabelsep}{0pt}
\defbibheading{pubtype}{\cvsubsection{#1}}
\renewcommand{\bibsetup}{\vspace*{-\baselineskip}}
\AtEveryBibitem{%
  \iffieldundef{doi}{}{\clearfield{url}}%
}


%% sample.bib contains your publications
\addbibresource{sample.bib}

\begin{document}
\name{Omar Cusma Fait}
\tagline{Data Science and Artificial Intelligence Student at the University of Trieste}
%% You can add multiple photos on the left or right

% profile picture
\photoR{2.7cm}{omar_profile_pic}

% \photoL{2.5cm}{Yacht_High,Suitcase_High}

\personalinfo{%
  % Not all of these are required!
  \email{omar.cusma.fait@gmail.com}
  \phone{+39 334 1819464}
  %\mailaddress{Åddrésş, Street, 00000 Cóuntry}
  \location{Friuli Venezia Giulia, ITALY}
  \homepage{https://omega97.github.io/}
  %\twitter{@twitterhandle}
  \linkedin{www.linkedin.com/in/omar-cusma-fait-27589a307}
  \github{Omega97}
  %\orcid{0000-0000-0000-0000}
  
  %% You can add your own arbitrary detail with
  %% \printinfo{symbol}{detail}[optional hyperlink prefix]
  % \printinfo{\faPaw}{Hey ho!}[https://example.com/]

  %% Or you can declare your own field with
  %% \NewInfoFiled{fieldname}{symbol}[optional hyperlink prefix] and use it:
  % \NewInfoField{gitlab}{\faGitlab}[https://gitlab.com/]
  % \gitlab{your_id}
  %%
  %% For services and platforms like Mastodon where there isn't a
  %% straightforward relation between the user ID/nickname and the hyperlink,
  %% you can use \printinfo directly e.g.
  % \printinfo{\faMastodon}{@username@instace}[https://instance.url/@username]
  %% But if you absolutely want to create new dedicated info fields for
  %% such platforms, then use \NewInfoField* with a star:
  % \NewInfoField*{mastodon}{\faMastodon}
  %% then you can use \mastodon, with TWO arguments where the 2nd argument is
  %% the full hyperlink.
  % \mastodon{@username@instance}{https://instance.url/@username}
}

\makecvheader
%% Depending on your tastes, you may want to make fonts of itemize environments slightly smaller
% \AtBeginEnvironment{itemize}{\small}

%% Set the left/right column width ratio to 6:4.
\columnratio{0.6}

% Start a 2-column paracol. Both the left and right columns will automatically
% break across pages if things get too long.
\begin{paracol}{2}


%%%%%%%%%%%%%%%%%%%%%%%%%%%%%%%%%%%%%%%%%%%%%%%%

\cvsection{About me}

I possess a strong background in \textbf{mathematics} and \textbf{physics}, which has refined my \textbf{analytical thinking} and \textbf{problem-solving} skills. These abilities have proven invaluable in approaching complex challenges from a logical perspective.

%I have developed familiarity with complex problems in \textbf{mathematics} and \textbf{physics}, my \textbf{analytical thinking}, and \textbf{problem-solving} skills that have proven valuable in approaching challenges from a logical perspective.
%My experience with chess taught me to stay \textbf{calm under pressure} and to seek a better solution to any given problem.
%I tend to dedicate myself fully to projects that I am passionate about.


%%%%%%%%%%%%%%%%%%%%%%%%%%%%%%%%%%%%%%%%%%%%%%%%

\cvsection{Work Experience}

\cvevent{Math Teacher}{ISIS Gregor\v{c}i\v{c}}{March 2025 -- June 2025}{Gorizia}
\begin{itemize}
\item Taught mathematics at the high school level, preparing and delivering lessons, supporting students in building solid problem-solving foundations, and assessing progress across the school term.
\end{itemize}


%%%%%%%%%%%%%%%%%%%%%%%%%%%%%%%%%%%%%%%%%%%%%%%%

%\cvsection{Things I am Proud Of}

%\cvachievement{\faStar}{My Personal Projects}{I apply the knowledge I acquired during the courses to create something...}

%\divider
%\cvachievement{\faSearch}{Scienze Under 18}{Co-presenter of “Scienze under 18”, an in-person activity that aims to bring math to the attention of the public and young students (supervisor: Caterina Vicentini)}

%\divider
%\cvachievement{\faTrophy}{Chess}{20 Chess awards, where I ranked 1st to 3rd in official tournament}


%%%%%%%%%%%%%%%%%%%%%%%%%%%%%%%%%%%%%%%%%%%%%%%%

\cvsection{Projects}

% Astrophysics Internship
\cvevent{Astrophysics Internship}{}{2023}{INAF} % 2023
\begin{itemize}
\item In a Study of the local correlation between two HEALPix images, I faced an open-ended problem for which I had to come up with my own solution to the problem.
\end{itemize}

\divider

% Semantic Retrieval System with LLMs
\cvevent{Semantic Retrieval System with LLMs}{}{}{} % {2024}{UniTS}
\begin{itemize}
\item I created a system that can search through a folder of documents to find the most relevant one, highlight the most relevant text in the document, or group the files in folders based on meaning. 
\end{itemize}

\divider

% Royal Game of Ur Self-Reinforcement Learning AI
\cvevent{Royal Game of Ur Self-Reinforcement Learning AI}{}{}{} %{2024}{UniTS}  
\begin{itemize}
\item I applied my knowledge in Reinforcement Learning by creating a bot that learns to play well at the game through self-play only.
\end{itemize}

\divider

% Wikipedia Recommendation System
\cvevent{Wikipedia Recommendation System}{}{}{} %{2024}{UniTS}
\begin{itemize}
%\items todo better
\item With my team, we created an online platform that recommends Wikipedia pages based on your preferences. I later improved the system by leveraging a lightweight sentence embedding model.
\end{itemize}

\divider

% Quantum Ground State Simulation
\cvevent{Quantum Ground State Simulation}{}{}{} %{2022}{UniTS} 
\begin{itemize}
\item I applied my knowledge in physics to write a program to find the ground state of a quantum system of interacting particles.
%\item Given the potential, the program will find the system's ground state using the variational method. 
%\item The program distinguishes between bosons and fermions and can also handle potentials with a singularity at the origin.
\end{itemize}

\divider

% Peak Detection
\cvevent{Peak Detection in Noisy Signal}{}{}{} % {2023}{}
\begin{itemize}
\item I helped a friend in automating the detection and localization of peaks in noisy spectroscopic data.
\end{itemize}

\divider

% Telegram Bot
\cvevent{Telegram Bot}{}{}{} % 2024
\begin{itemize}
\item I created a general-purpose bot that manages and interacts with Telegram users in real time.
\end{itemize}


%\divider

% Wave simulation
%\cvevent{Wave simulation}{UniTS}{}{} % 2022
%\begin{itemize}
%\item The library implements a simple simulation of a 2D wave using the Runge-Kutta method that can be exported to a GIF.
%\item The library allows to specify the boundary conditions, initial conditions, and obstacles. 
%\end{itemize}
%\divider

% Google Hash 2019
%\cvevent{Google Hash 2019}{}{}{} % 2019
%\begin{itemize}
%\item I had the opportunity to show my programming and teamwork skills in an intense competition
%\end{itemize}
%\divider

%%%%%%%%%%%%%%%%%%%%%%%%%%%%%%%%%%%%%%%%%%%%%%%%

\cvsection{Interests}

\cvtag{Python} 
\cvtag{Mathematics} 
\cvtag{Physics} 
\cvtag{Science}\\
\cvtag{Simulations} 
\cvtag{Reinfocement Learning} 
\cvtag{AI} \\
\cvtag{Machine Learning} 
\cvtag{Big Data}
\cvtag{Deep Learning} 
 

%%%%%%%%%%%%%%%%%%%%%%%%%%%%%%%%%%%%%%%%%%%%%%%%
\switchcolumn

\cvsection{Digital Skills}

\cvskill{Python}{4.5}
\divider

\cvtag{Pandas}
\cvtag{NumPy}
\cvtag{MatPlotLib}
\cvtag{GitHub}\\
\cvtag{PyCharm}
\cvtag{Latex}
\cvtag{Krita}
\cvtag{Windows}\\
\cvtag{Excel}
\cvtag{Word}
\cvtag{Pascal}
\cvtag{PowerPoint}


%%%%%%%%%%%%%%%%%%%%%%%%%%%%%%%%%%%%%%%%%%%%%%%%

\cvsection{Languages}

\cvskill{Italian (native)}{5}
\divider

\cvskill{Slovenian (native)}{4.5}
\divider

\cvskill{English (B2)}{4.5} %% Supports X.5 values.
\divider

Learning: \cvtag{Japanese}

\medskip


%%%%%%%%%%%%%%%%%%%%%%%%%%%%%%%%%%%%%%%%%%%%%%%%

\cvsection{Education}

\cvevent{Currently pursuing: Master's Degree\ in Data Science and Artificial Intelligence}{Università degli Studi di Trieste}{}{}

\divider
\cvevent{Bachelor's Degree\ in Physics}{Università degli Studi di Trieste}{}{}
Thesis title: A Comparative Analysis of Neural Network Methonds for Source Detection

\divider
\cvevent{High School Diploma\ Applied Sciences}{Istituto di Istruzione Superiore Gabriele D'Annunzio}{}{}

%\divider
%\cvevent{M.Sc.\ in Your Discipline}{Your University}{Sept 2001 -- June 2002}{}

%\divider
%\cvevent{B.Sc.\ in Your Discipline}{Stanford University}{Sept 1998 -- June 2001}{}

% \divider


%%%%%%%%%%%%%%%%%%%%%%%%%%%%%%%%%%%%%%%%%%%%%%%%

\cvsection{Referees}

% \cvref{name}{email}{mailing address}
\cvref{Dr.\ Elena Orlando}{INAF}{orland.ele@gmail.com}

\divider

\cvref{Peter \v{C}ernic}{Principal, ISIS Gregor\v{c}i\v{c} -- Gorizia}{Available upon request}

%%%%%%%%%%%%%%%%%%%%%%%%%%%%%%%%%%%%%%%%%%%%%%%%

\end{paracol}

\end{document}
